% =====================
% Chapter 1: Python Basics
% Audience: Architecture students new to coding
% Sources aligned with: A Byte of Python (About Python, Installation, First Steps)
% =====================

\h{Getting Started With Python}

\begin{NoteBox}
\textbf{What you'll learn.} What Python is and why we use it; how to install Python 3 and set up \emph{Visual Studio Code}; how to run a script and use the interactive prompt; and the first building blocks (print, variables, comments, indentation). Examples use unit-aware names (e.g., \texttt{span\_m}) tied to studio or structural tasks.
\end{NoteBox}

\hh{Why Python (for architects and engineers)}
Python is a high-level, interpreted, multi-paradigm language designed for readability:
\begin{itemize}
  \item Readable, concise syntax (indentation instead of braces).
  \item ``Batteries included'': rich standard library (files, text, data, math, web).
  \item Cross-platform and free (Windows, macOS, Linux; open source).
  \item Versatile: quick scripts, data wrangling (CSV/Excel), visualization, geometry, APIs, automation.
\end{itemize}
For ARC500, Python helps you prototype: check areas, generate reports, process schedules, or test structural formulas without heavy software.

\hh{Interpreted and interactive}
\begin{itemize}
  \item \textbf{Interpreted}: run code directly (no separate compile step) for rapid feedback.
  \item \textbf{Interactive}: a prompt (REPL) lets you try expressions and see results immediately.
\end{itemize}

\hh{Installation and setup (Python 3 + VS Code)}
\textbf{Goal:} Have Python 3 installed, VS Code configured, and a working terminal.

\hhh*{1) Install Python 3}
\begin{itemize}
  \item \textbf{Windows}: \url{https://www.python.org/downloads/}. During setup, check ``Add Python to PATH''. If prompted, accept ``Disable PATH length limit''.
  \item \textbf{macOS}: Use the official installer or a package manager. You will typically use \texttt{python3} in Terminal.
\end{itemize}

\code{c01-verify-python}{bash}{
    python --version
    python3 --version
    py --version
}{Verify Python in your terminal}

\hhh*{2) Install Visual Studio Code}
\begin{itemize}
  \item Install VS Code: \url{https://code.visualstudio.com/}.
  \item Install the \textbf{Python} extension (by Microsoft).
  \item Press \texttt{Ctrl+Shift+P}/\texttt{Cmd+Shift+P} and run ``Python: Select Interpreter'' to choose Python 3.
\end{itemize}

\hhh*{3) Open a folder and the integrated terminal}
\begin{itemize}
  \item File \textrightarrow{} Open Folder... (e.g., \texttt{arc500/}).
  \item Terminal \textrightarrow{} New Terminal to open a shell in that folder.
\end{itemize}

\hh{Your first script in VS Code}
Create \texttt{hello.py} and paste:

\code{c01-hello}{python}{
    print("Hello, ARC500!")
    for i in range(3):
        print("Iteration", i)
}{First run in Python}

Run from the VS Code terminal:

\code{c01-run}{bash}{
    python hello.py    # Windows or macOS if python -> Python 3
    py hello.py        # Windows launcher alternative
    python3 hello.py   # macOS typical
}{Run your script from the terminal}

You should see the greeting and lines counting 0, 1, 2.

\hhh{Troubleshooting quick fixes}
\begin{itemize}
  \item \textbf{Command not found}: Python not on PATH. Reinstall and check ``Add Python to PATH'' (Windows). On macOS try \texttt{python3}.
  \item \textbf{Wrong interpreter}: Use ``Python: Select Interpreter'' in VS Code.
  \item \textbf{Nothing prints}: Be sure you ran the file from the terminal, not just typed in the editor.
\end{itemize}

\hh{Using the interactive prompt (REPL)}
Start it with \texttt{python}, \texttt{python3}, or \texttt{py} and try:

\code{c01-repl}{python}{
    >>> 2 + 3 * 4
    14
    >>> width_m, length_m = 7.2, 9.1
    >>> area_m2 = width_m * length_m
    >>> area_m2
    65.52
    >>> f"Area = {area_m2:.1f} m^2"
    'Area = 65.5 m^2'
}{Quick calculations in the REPL}

Exit with \texttt{exit()} or Ctrl+D (macOS/Linux) or Ctrl+Z then Enter (Windows).

\hh{First steps: print, variables, and input}
\texttt{print} shows results; variables store values; \texttt{input()} reads user text (always a string).

\code{c01-io}{python}{
    print("Room calculator")
    w = float(input("Width (m): "))
    l = float(input("Length (m): "))
    area_m2 = w * l
    print(f"Area: {area_m2:.2f} m^2")
}{Input, conversion, and f-string output}

\begin{NoteBox}
\textbf{Key idea.} \texttt{input()} returns a string. Convert with \texttt{float()} or \texttt{int()} before math. Keep unit-aware names (e.g., \texttt{height\_m}, \texttt{udl\_kPa}).
\end{NoteBox}

\hh{Comments and readable code}
Comments (\texttt{\#}) do not execute; they document intent, assumptions, and nonobvious choices.

\code{c01-comments}{python}{
    # Compute total facade area for a rectangular elevation
    width_m = 18.0    # meters
    height_m = 12.0   # meters
    facade_area_m2 = width_m * height_m
    print(f"Facade area: {facade_area_m2} m^2")
}{Comments and unit-aware names}

\hh{Indentation defines blocks}
Python uses significant indentation to group code. Use 4 spaces per level (no tabs).

\code{c01-indent}{python}{
    threshold = 3
    for i in range(5):
        if i >= threshold:
            print(f"{i} is at or above threshold")
        else:
            print(f"{i} is below threshold")
}{Blocks defined by indentation}

\begin{NoteBox}
\textbf{Common errors.} \emph{IndentationError}: misaligned spaces or mixing tabs/spaces. \emph{NameError}: using a name before assignment. Read the traceback from the bottom: the last line names the error type and message.
\end{NoteBox}

\hh{Numbers and simple expressions (first look)}
Use Python like a calculator; include units in names.

\code{c01-exprs}{python}{
    span_m = 8.0
    uniform_load_kNpm = 12.5
    midspan_M_kNm = uniform_load_kNpm * span_m**2 / 8
    print(f"M_mid ~= {midspan_M_kNm:.1f} kN*m")
}{Expressions with units and exponentiation}

\hh{Strings and f-strings (first look)}
F-strings embed values and control formatting:

\code{c01-fstrings}{python}{
    project = "Studio Pavilion"
    level = 2
    print(f"{project} - Level {level}")
    print(f"{1/3:.3f}")       # 0.333
    print(f"{'Lobby':_^11}")  # ___Lobby___
}{Readable formatted output}

\hh{Mini-Lab 1 -- Room Area and Volume (15 min)}
\textbf{Goal.} Prompt for width, length, and height (m); compute area and volume; print a one-line summary.
\begin{itemize}
  \item Use \texttt{input()} then \texttt{float()} for numbers.
  \item Use unit-aware names.
  \item Use an f-string with controlled decimals.
\end{itemize}

\code{ml01-starter}{python}{
    print("Room calculator")
    w = float(input("Width (m): "))
    l = float(input("Length (m): "))
    h = float(input("Height (m): "))
    area_m2 = w * l
    volume_m3 = area_m2 * h
    print(f"Area = {area_m2:.2f} m^2, Volume = {volume_m3:.2f} m^3")
}{Mini-Lab starter}

\textbf{Challenge.} Ask for a material unit weight (kN/m\textsuperscript{3}) and estimate total weight of the room's air volume.

\hh{Style tips (PEP 8 essentials)}
\begin{itemize}
  \item Four spaces per indent; keep lines under about 79--100 characters.
  \item \texttt{lowercase\_with\_underscores} for variables and functions.
  \item Meaningful, unit-aware names: \texttt{span\_m}, \texttt{live\_load\_psf}.
\end{itemize}

\hh{Quick practice set}
Create \texttt{basics\_drills.py} and solve:
\begin{enumerate}
  \item Ask for total floor area (m\textsuperscript{2}) and floor-to-floor height (m). Print an estimate of exterior wall area using a simple rule you state in a comment (e.g., perimeter \textasciitilde{} 0.35*sqrt(area)).
  \item Print the first 10 even integers on one line, separated by commas.
  \item Given \texttt{span\_m} and \texttt{uniform\_load\_kNpm}, compute midspan moment \(wL^2/8\) (kN*m) and print it with units using an f-string.
\end{enumerate}
